\section{Diagnóstico de la Estructura de Gestión del Suelo}

A partir de la base de datos espacial normalizada del Plan General de Ordenación (formato FIP 2.0 \cite{sitcan_planeamiento_2026}
  % integrado en PostGIS
con ayuda del entorno de computación agéntica Google Antigravity \cite{google_antigravity_2026}),
se analizó la estructura de gestión del suelo en el municipio.

\subsection{Distribución de las Unidades de Actuación (UAs) por Categorías del Suelo}

El análisis de las UAs delimitadas en el término municipal refleja una concentración determinante en el suelo urbano sujeto a procesos de regularización y renovación:

\begin{table}[htbp]
  \centering
  \caption{Resumen de Unidades de Actuación por Categoría del Suelo}
  \label{tab:resumen_suelo}
  \vspace{0.2cm}
  \begin{tabular}{lrrr}
  \toprule
  \textbf{Categoría del Suelo} & \textbf{Nº UAs} & \textbf{Edificabilidad ($m^2c$)} & \textbf{Superficie ($m^2$)} \\
  \midrule
  Suelo Urbano No Consolidado & 125 & 1.557.793 & 2.915.364 \\
  Suelo Urbanizable Sectorizado & 3 & 80.778 & 112.021 \\
  \midrule
  \textbf{Total} & \textbf{128} & \textbf{1.638.571} & \textbf{3.027.385} \\
  \bottomrule
\end{tabular}

\end{table}

De las 128 unidades delimitadas, el \textbf{97,7\%} (125 UAs) corresponden a \textbf{Suelo Urbano No Consolidado (SUNC)}, sumando una capacidad edificatoria lucrativa de más de 1,55 millones de $m^2c$. Por su parte, el \textbf{Suelo Urbanizable Sectorizado (SUS)} concentra únicamente 3 unidades (ubicadas en el ámbito residencial de El Tablero), con 80.778 $m^2c$.

\begin{figure}[htbp]
  \centering
  \includegraphics[width=0.92\textwidth]{figures/chart_land_classes.pdf}
  \caption{Capacidad edificatoria y número de UAs por categoría de suelo.}
  \label{fig:chart_suelo}
\end{figure}

\subsection{Sistemas de Ejecución y Gestión Urbanística}

El régimen de desarrollo de las actuaciones delimitadas determina el grado de iniciativa privada frente a la intervención directa de la administración municipal, conforme al marco legal autonómico regulado en el Título V de la Ley 4/2017, del Suelo y de los Espacios Protegidos de Canarias \cite{ley_suelo_canarias_2017}:

\begin{table}[htbp]
  \centering
  \caption{Distribución de Unidades de Actuación según Sistema de Ejecución}
  \label{tab:resumen_sistemas}
  \vspace{0.2cm}
  \begin{tabular}{p{6.8cm}rrr}
  \toprule
  \textbf{Sistema de Ejecución} & \textbf{Nº UAs} & \textbf{\% UAs} & \textbf{Edificabilidad ($m^2c$)} \\
  \midrule
  Sistemas de ejecución privada legalmente establecidos & 75 & 58.6\% & 1.169.062 \\
  Cooperación & 45 & 35.2\% & 463.292 \\
  No especificado & 3 & 2.3\% & 6.217 \\
  Expropiación & 3 & 2.3\% & - \\
  Convenio urbanístico & 2 & 1.6\% & - \\
  \midrule
  \textbf{Total} & \textbf{128} & \textbf{100.0\%} & \textbf{1.638.571} \\
  \bottomrule
\end{tabular}

\end{table}

\begin{itemize}
  \item \textbf{Compensación / Ejecución Privada (58,6\%)}: Mecanismo predominante, acumulando más de 1,16 millones de $m^2c$ edificables donde los propietarios asumen la redacción de los proyectos y las cesiones obligatorias (regulado en el \href{https://www.boe.es/buscar/act.php?id=BOE-A-2017-10295#ci-21}{Capítulo III} \cite{ley_suelo_canarias_2017}
      %; consulta en \href{run:supporting_docs/BOE-A-2017-10295-consolidado.pdf#page=115}{PDF local, pág. 115}
    ).
  \item \textbf{Cooperación Pública (35,2\%)}: Abarca 45 unidades (463.292 $m^2c$), donde la ejecución de las obras corresponde al Ayuntamiento a costa de los propietarios (\href{https://www.boe.es/buscar/act.php?id=BOE-A-2017-10295#ci-22}{Capítulo IV}
      %; consulta en \href{run:supporting_docs/BOE-A-2017-10295-consolidado.pdf#page=127}{PDF local, pág. 127}
    ).
  \item \textbf{Expropiación y Convenios (3,9\%)}: Reservados para ámbitos con afecciones de sistemas generales o actuaciones asistemáticas singulares.
\end{itemize}

\subsection{Planificación y Realidad: Verificación Aerofotogramétrica}

El modelo de datos del FIP 2.0 constituye una proyección normativa y temporal de lo que debía ejecutarse en un horizonte legal máximo de 8 años (cuatrienios 2015--2019 y 2019--2023). Al haber expirado dicho plazo legal, la contrastación espacial entre las poligonales de las Unidades de Actuación (\texttt{v\_unidades\_actuacion}) y la cobertura aerofotogramétrica ortorrectificada de máxima actualidad de IDECanarias/GRAFCAN (\emph{OrtoExpress Urbana}, 10--12,5 cm/píxel) permite auditar de forma empírica el grado real de transformación del suelo sobre el territorio.

La inspección sobre los principales nodos de desarrollo municipal revela un patrón sistemático de inejecución y parálisis:

\subsubsection{Suelo Urbanizable Sectorizado: El Tablero (SUS 5.13)}

\begin{figure}[htbp]
  \centering
  \includegraphics[width=0.88\textwidth]{figures/orto_detalle_el_tablero.jpg}
  \caption{Superposición de las UAs del Sector SUS 5.13 (El Tablero) sobre la ortofoto aérea reciente de IDECanarias. Fuente: Elaboración propia a partir de datos del FIP 2.0 e IDECanarias.}
  \label{fig:orto_el_tablero}
\end{figure}

Las tres unidades delimitadas en el sector (\textbf{UA-1, UA-2 y UA-3 de SSU 5.13},
junto con la contigua \textbf{UA-RT-1 del Ámbito 5.16.2}) concentran la totalidad del suelo urbanizable residencial sectorizado del municipio,
con más de 80.000 $m^2c$ de edificabilidad lucrativa programada.
Como se aprecia en la ortofoto (Figura~\ref{fig:orto_el_tablero})
la totalidad del ámbito permanece como una gran explanada terrosa y baldía encajada entre el núcleo preexistente y el campo de fútbol.
No se ha ejecutado la red viaria estructurante, exceptuando la \textbf{UA-1}, no existen aceras recepcionadas ni se ha edificado una sola de las viviendas previstas.

\subsubsection{Suelo Urbano No Consolidado: Cabo Llanos (Ámbito 2.5.3)}

\begin{figure}[htbp]
  \centering
  \includegraphics[width=0.88\textwidth]{figures/orto_detalle_cabo_llanos.jpg}
  \caption{Unidad de Actuación UA-RT-1 en el Ámbito 2.5.3 (Cabo Llanos) sobre ortofoto aérea. Fuente: Elaboración propia a partir de datos del FIP 2.0 e IDECanarias.}
  \label{fig:orto_cabo_llanos}
\end{figure}

La \textbf{UA-RT-1 del Ámbito 2.5.3} acumula una capacidad lucrativa superior a los 25.000 $m^2c$ en el área de mayor valor terciario y residencial de la capital. La imagen aérea (Figura~\ref{fig:orto_cabo_llanos}) muestra que el polígono, delimitado en rojo y rodeado por manzanas residenciales de alta densidad completamente colmatadas, constituye un vacío urbano absoluto: un solar de tierra y polvo utilizado informalmente como descampado, sin que se haya materializado la reparcelación ni el proyecto de urbanización.

\subsubsection{Suelo Urbano No Consolidado: Barranco de Santos (La Ladera)}

\begin{figure}[htbp]
  \centering
  \includegraphics[width=0.72\textwidth]{figures/orto_detalle_barranco_santos.jpg}
  \caption{Unidades de Actuación en La Ladera (Barranco de Santos) sobre ortofoto aérea. Fuente: Elaboración propia a partir de datos del FIP 2.0 e IDECanarias.}
  \label{fig:orto_barranco_santos}
\end{figure}

El frente urbano sobre el Barranco de Santos refleja la complejidad del tejido intersticial (Figura~\ref{fig:orto_barranco_santos}):
\begin{itemize}
  \item \textbf{Ámbito 1.7.4A1 (La Ladera A1)}: Presenta un mosaico degradado de edificaciones irregulares preexistentes y una explanada terrosa empleada como estacionamiento no regulado.
  \item \textbf{Ámbito 1.7.4A2 (La Ladera A2)}: Es una ladera escarpada de vegetación seca que vierte directamente al cauce, carente de viales, encintado de aceras o infraestructuras básicas de servicios.
\end{itemize}

\clearpage
\subsubsection{Suelo Urbano No Consolidado: Montaña de Taco}

\begin{figure}[htbp]
  \centering
  \includegraphics[width=0.72\textwidth]{figures/orto_detalle_montana_taco.jpg}
  \caption{Ámbitos de gestión en Montaña de Taco (Industrial, Parque y Residencial) sobre ortofoto aérea. Fuente: Elaboración propia a partir de datos del FIP 2.0 e IDECanarias.}
  \label{fig:orto_montana_taco}
\end{figure}

En el sector limítrofe de Montaña de Taco confluyen tres unidades con distinta vocación funcional (Figura~\ref{fig:orto_montana_taco}):
\begin{itemize}
  \item \textbf{Ámbito 4.9.4 (Parque Montaña de Taco)}: Permanece como una gran cantera de extracción de áridos a cielo abierto, sin haber acometido la restauración ambiental ni la ejecución del sistema general de parque urbano.
  \item \textbf{Ámbito 4.9.5 (Montaña de Taco Residencial)}: Mantiene su condición de ladera baldía y terrosa sin traza viaria ni urbanización.
  \item \textbf{Ámbito 4.9.3 (Montaña de Taco Industrial)}: Agrupa naves fabriles obsoletas y patios de acopio industrial sin el proceso de renovación urbana exigido por el planeamiento.
\end{itemize}

Este contraste evidencia que la inoperancia de los sistemas de ejecución privada (compensación) y la parálisis administrativa derivada de la anulación judicial del PGO de 2013/2015 han mantenido bloqueada la bolsa de suelo lucrativo del municipio, impidiendo que la oferta de suelo programada diera respuesta a las necesidades habitacionales y funcionales de Santa Cruz de Tenerife.
